% Section: geometric decoupling of an audit probe (draft v3).
% v3 = v2 + corrections from the independent counter-reading (papiers/D/audits/contre-lecture-decouplage-geometrique-v1/):
%      major M1-M4 and minors m1-m11. Written by the second daughter session (Opus), NOT by Fable.
% v3-correctif-1: single correction loop after the fidelity check of v3 (verification-decouplage-v3-v1.md, defects 1-12).
% STATUS: A REVOIR PAR FABLE (piece de fond); not for diffusion.
% Citation keys (all TO BE VERIFIED ON PDF, rule R7; see decouplage-geometrique-v3-references.md):
%   pearson1903, aitken1934, lawley1943, stein1981, hirschfeld1935, gebelein1941, renyi1959,
%   lancaster1957, kolmogorov1960, chapman1951, jacot2018, cloud2024, bailey2024, gupta2025.
% Macros expected from the host paper: \E, \Cov, \Var; label thm:selection (static selection theorem).
% Numerical companion: code/experience_jouet_decouplage_sonde_v1.py (official figures);
%   counter-examples: audits/contre-lecture-decouplage-geometrique-v1/verif/v2_jouet.py, v3_moniteur_et_signe.py,
%   v4_contre_exemple_signe.py, all replayed by the daughter session on a copy (v2_jouet.py replayed for v3-correctif-1).

\section{Geometry of Audit Decoupling: Probes on Internal Representations}\label{sec:geometry}

The static decoupling regime (R2) is stated for abstract survival probabilities. We now make it concrete for the case that motivates it in practice: an in-loop filter that reads a chain-of-thought (text) head of a network, and an audit probe that reads internal activations of the same network and is not itself used as a training signal. Keeping the probe out of the loop is not what protects it: in the static model the relevant quantity is the dependence between in-loop survival and audit survival (an exact covariance under rejection; independence or negative association under the penalty, non-positive covariance being necessary but not sufficient). The two results below play different roles. Theorem~\ref{thm:geometry-selection} characterizes that dependence under \emph{selection}: when the law of the surviving episodes is reweighted by the filter, which probes see their score law unchanged. Theorem~\ref{thm:geometry-training} concerns only the \emph{instrument}: when the representation is trained against the filter, which probes keep the same score function $x \mapsto s(x)$. An unchanged instrument does not imply an unchanged audit among the episodes that survive the filter (Remark~\ref{rem:instrument-vs-audit}). Both results apply well-known tools---selection formulas under normality, Stein's identity, maximal correlation, the neural tangent kernel and the masking of gradients as in gradient routing---to the filter-versus-audit setting; we claim no novelty for the tools themselves.

\paragraph{Setting.} An episode has an internal representation $h \in \mathbb{R}^d$. The in-loop filter only sees the readout $z = W h \in \mathbb{R}^m$ of a chain-of-thought head $W \in \mathbb{R}^{m \times d}$. The audit is a frozen linear probe with score $s = w^\top h$, $w \in \mathbb{R}^d \setminus \{0\}$, oriented so that larger scores indicate cheating. We study two ways in which in-loop filtering acts on the population of cheating episodes: (i) \emph{selection}, where the law $P$ of $h$ on a cheating stratum is reweighted by a function of $z$ (rejection sampling, best-of-$N$, Gibbs tilting), and (ii) \emph{gradient training} of the representation itself against a loss that depends on $h$ only through the heads.

\subsection{Selection on the filter readout}

For a measurable $g : \mathbb{R}^m \to [0, \infty)$ with $0 < \E_P[g(Wh)] < \infty$, write $P_g(dh) := g(Wh)\, P(dh) / \E_P[g(Wh)]$ for the selected law. Operator~A on an iso-reward stratum is the case $g = \beta_{\mathrm{loop}}$, the survival probability of the filter.

\begin{theorem}[Selection on a text readout]\label{thm:geometry-selection}
Let $h \sim P$ with $\E_P\|h\|^2 < \infty$.
\begin{enumerate}
  \item \textbf{Exact invariance.} The law of $s = w^\top h$ under $P_g$ equals its law under $P$ for every admissible $g$ if and only if $s$ and $Wh$ are independent under $P$. If $P = \mathcal{N}(\mu, \Sigma)$, this holds if and only if $W \Sigma w = 0$.
  \item \textbf{Exact mean shift (Gaussian case).} If $P = \mathcal{N}(\mu, \Sigma)$ and $g > 0$ is $C^1$ with $\E_P[g(Wh)] < \infty$ and $\E_P\|\nabla g(Wh)\| < \infty$, then
  \begin{equation}\label{eq:stein-shift}
    \E_{P_g}[s] - \E_P[s] = (W \Sigma w)^\top\, \E_{P_g}\!\big[\nabla \log g(Wh)\big].
  \end{equation}
  For a logistic filter that intercepts with probability $\sigma(a^\top z - b)$ and survival weight $g(z) = 1 - \sigma(a^\top z - b)$,
  \begin{equation}\label{eq:logistic-shift}
    \E_{P_g}[s] - \E_P[s] = -\,\E_{P_g}\!\big[\sigma(a^\top W h - b)\big]\; \langle W^\top a, w \rangle_{\Sigma},
    \qquad \langle u, v \rangle_\Sigma := u^\top \Sigma v.
  \end{equation}
  Hence selection moves the mean probe score of surviving cheaters towards the honest side if $\langle W^\top a, w\rangle_\Sigma > 0$, away from it if $\langle W^\top a, w\rangle_\Sigma < 0$, and leaves the whole probe-score law unchanged if $W\Sigma w = 0$.
  \item \textbf{Robustness under partial coupling (any law).} Assume $\Var_P(s) > 0$ (otherwise $s$ is a.s.\ constant and there is nothing to bound). Let $\rho^\star := \sup\{\operatorname{Corr}_P(\phi(s), \psi(Wh))\}$ be the Hirschfeld--Gebelein--R\'enyi maximal correlation~\cite{hirschfeld1935,gebelein1941,renyi1959} between $s$ and $Wh$ under $P$ (supremum over non-constant square-integrable $\phi, \psi$). If $\E_P[g(Wh)^2] < \infty$, then for every threshold $t$,
  \begin{align}
    \big|\E_{P_g}[s] - \E_P[s]\big| &\le \rho^\star\, \operatorname{sd}_P(s)\, \sqrt{\chi^2(P_g \,\|\, P)}, \label{eq:robust-mean}\\
    \big|P_g(s > t) - P(s > t)\big| &\le \min\Big\{1,\ \rho^\star \sqrt{P(s>t)\,(1 - P(s>t))\, \chi^2(P_g \,\|\, P)}\Big\} \notag\\ &\le \min\Big\{1,\ \tfrac12 \rho^\star \sqrt{\chi^2(P_g \,\|\, P)}\Big\}, \label{eq:robust-tpr}
  \end{align}
  where $\chi^2(P_g \| P) = \Var_P(g(Wh)) / \E_P[g(Wh)]^2$. For Operator~A ($g = \beta_{\mathrm{loop}}$) this is $\Var(\beta_{\mathrm{loop}})/\E[\beta_{\mathrm{loop}}]^2$, the same variance that drives conditional enrichment in Theorem~\ref{thm:selection}. For Gaussian $P$, $\rho^\star$ equals the canonical correlation $\big(w^\top \Sigma W^\top (W\Sigma W^\top)^{+} W \Sigma w / w^\top \Sigma w\big)^{1/2}$~\cite{lancaster1957,kolmogorov1960}. Both bounds remain true, and are tighter, with $\rho^\star$ replaced by the maximal correlation between $s$ and $g(Wh)$, which is at most $\rho^\star$.
\end{enumerate}
\end{theorem}

Part~2 and the Gaussian case of Part~1 follow from the classical selection formulas under multivariate normality~\cite{pearson1903,aitken1934,lawley1943}, here written through Stein's identity~\cite{stein1981}; the general case of Part~1 is an elementary fact about indicator weights; Part~3 is a Cauchy--Schwarz bound of Chapman--Robbins type~\cite{chapman1951} combined with maximal correlation. The first inequality in~\eqref{eq:robust-tpr} is attained: for $s = z$, $g = \mathbf{1}\{z \le t\}$ and $\rho^\star = 1$ the change equals the bound. Note that the orthogonality relevant to selection is $W\Sigma w = 0$ (orthogonality in the metric $\Sigma$), not $Ww = 0$; the two coincide for every $w$ when the row space of $W$ is invariant under $\Sigma$ (for instance $\Sigma \propto I$, or $\Sigma = \operatorname{diag}(1,2,3)$ with $W = e_1^\top$), and differ in general. Theorem~\ref{thm:geometry-training}(2) uses $Ww = 0$, so a probe that is invariant under gradient training of a linear network need not be invariant under selection.

\begin{proof}
\emph{Part 1.} If $s \perp Wh$, then for bounded measurable $\phi$, $\E_P[\phi(s) g(Wh)] = \E_P[\phi(s)]\,\E_P[g(Wh)]$, so the law of $s$ under $P_g$ is that under $P$. Conversely, for any Borel $B \subset \mathbb{R}^m$ with $P(Wh \in B) > 0$, $g = \mathbf{1}_B$ is admissible and invariance gives $P(s \in A, Wh \in B) = P(s \in A) P(Wh \in B)$ for all Borel $A$; the identity is trivial when $P(Wh \in B) = 0$; hence independence. In the Gaussian case $(s, Wh)$ is jointly Gaussian with cross-covariance $W \Sigma w$, and joint Gaussians are independent iff uncorrelated.

\emph{Part 2.} Stein's identity for $h \sim \mathcal{N}(\mu, \Sigma)$ ($\Sigma$ possibly singular) and $\psi(h) = g(Wh)$ with $\E\|\nabla\psi\| < \infty$ gives $\E[(h - \mu)\psi(h)] = \Sigma\, \E[\nabla \psi(h)] = \Sigma W^\top \E[\nabla g(Wh)]$. Taking the inner product with $w$ and dividing by $\E_P[g(Wh)] > 0$ yields $\E_{P_g}[s] - w^\top \mu = (W\Sigma w)^\top \E_P[\nabla g(Wh)]/\E_P[g(Wh)]$, and $\E_P[\nabla g]/\E_P[g] = \E_{P_g}[\nabla g / g]$ since $g > 0$. For the logistic survival weight, $\nabla \log(1 - \sigma(a^\top z - b)) = -\sigma(a^\top z - b)\, a$, and $(W \Sigma w)^\top a = \langle W^\top a, w\rangle_\Sigma$.

\emph{Part 3.} For any square-integrable $\phi$, $\E_{P_g}[\phi(s)] - \E_P[\phi(s)] = \Cov_P(\phi(s), g(Wh)) / \E_P[g(Wh)]$. By definition of $\rho^\star$, $|\Cov_P(\phi(s), g(Wh))| \le \rho^\star \operatorname{sd}_P(\phi(s))\operatorname{sd}_P(g(Wh))$ (trivially if either is constant). Divide by $\E_P[g]$ and use $\operatorname{sd}_P(g)/\E_P[g] = \sqrt{\chi^2(P_g\|P)}$, which holds because $dP_g/dP = g/\E_P[g]$. Take $\phi(s) = s$ for~\eqref{eq:robust-mean} and $\phi(s) = \mathbf{1}\{s > t\}$ for~\eqref{eq:robust-tpr}; the last inequality uses $x(1-x) \le 1/4$, and a difference of probabilities never exceeds $1$. The same argument with $\psi = $ identity on $g(Wh)$ gives the tighter version. For jointly Gaussian $(s, Wh)$ the maximal correlation equals the canonical correlation, which for a scalar $s$ is the stated multiple correlation coefficient.
\end{proof}

\begin{remark}[Uncorrelated is not enough]\label{rem:uncorrelated}
Outside the Gaussian case, $\Cov_P(s, Wh) = 0$ does not imply invariance. The example fits the linear-probe setting: $x \sim \mathcal{N}(0,1)$, $h = (x, x^2 - 1)$, $W = (1, 0)$, $w = (0, 1)$, so $z = x$ and $s = x^2 - 1$, with $\Sigma = \operatorname{diag}(1, 2)$ and $W\Sigma w = 0$. With the logistic survival weight $g(z) = 1 - \sigma(2z - 1/2)$, selection shifts the mean of $s$ by $-0.0775$ (exact, by one-dimensional quadrature; the companion simulation gives $-0.078$, standard error $5 \times 10^{-4}$). The relevant quantity is the maximal correlation $\rho^\star$ of Part~3, not the linear correlation.
\end{remark}

\subsection{Gradient training of a shared representation}

Selection reweights episodes without changing the network; gradient training changes the representation itself. Consider a one-hidden-layer network $h(x) = \varphi(Mx)$, $M \in \mathbb{R}^{d \times p}$ with rows $m_i^\top$, $\varphi$ applied coordinatewise, either $C^1$ or the ReLU (with the convention $\varphi'(0) = 0$). The training objective depends on $h$ only through the chain-of-thought head $W$ (the filter path) and a task head $U \in \mathbb{R}^{d \times k}$ (the reward path), plus neuron-separable anchoring terms on activations and on parameters:
\begin{equation}\label{eq:train-objective}
  \mathcal{L}(M) = \E_{x \sim \mathcal{D}}\big[\ell_x\big(W h(x),\, U^\top h(x)\big)\big] + \sum_{i=1}^d \E_{x \sim \mathcal{D}}\big[\psi_i\big(h_i(x), x\big)\big] + \sum_{i=1}^d r_i(m_i),
\end{equation}
where $\ell_x$, $\psi_i(\cdot, x)$ and $r_i$ are $C^1$, and each anchoring term is stationary at initialisation: $\partial_1 \psi_i(h_i^{(0)}(x), x) = 0$ and $\nabla r_i(m_i(0)) = 0$ (e.g.\ $\psi_i = \nu\, c(x) (h_i - h_i^{(0)})^2$ with a weight $c(x) \ge 0$ common to all neurons, anchoring to the pretrained representation, and $r_i = \lambda \|m_i - m_i(0)\|^2$, weight decay centred at the initial weights). Write $S := \{i : w_i \neq 0\}$ for the neurons read by the probe. Training is any first-order method, started from a zero optimizer state (zero momentum buffer, zero Adam moments), whose update of row $m_i$ may depend on all past gradients but vanishes whenever all past gradients of $\mathcal{L}$ with respect to $m_i$ vanish. This covers gradient descent with any step sizes, mini-batches drawn from $\mathcal{D}$, momentum, Adam with $\varepsilon > 0$ (with $\varepsilon = 0$ the update is $0/0$), per-coordinate clipping, global-norm clipping and layer-wise trust ratios without decay. It excludes full preconditioners (e.g.\ full natural gradient), optimizer states carried over from earlier training (a non-zero initial Adam first moment of $10^{-3}$ moves the probe's neurons, by $7 \times 10^{-7}$ in a test on the companion network), and plain or decoupled weight decay, which moves every neuron, the probe's included. (Absent activation anchoring on the probe's neurons and with $1 - \gamma\lambda > 0$, a pure decoupled decay multiplies the probe's rows by a common factor; for a positively homogeneous $\varphi$ this rescales scores and so changes a fixed threshold, but not the area under the ROC curve. With activation anchoring this is no longer a pure rescaling.)

\begin{theorem}[Neuron-level invariance of the instrument under gradient training]\label{thm:geometry-training}
\leavevmode
\begin{enumerate}
  \item \textbf{Neuron-disjoint probes keep their score function.} If $W_{:,i} = 0$ and $U_{i,:} = 0$ for all $i \in S$, then $m_i(t) = m_i(0)$ for all $i \in S$ and all $t$, so the score function $x \mapsto s(x) = w^\top h(x)$ is unchanged along the entire trajectory, whatever the filter loss, the reward loss, the step sizes and the number of steps. The statement is pointwise in $x$, so it also holds when the training inputs depend on the policy. It says nothing about the \emph{law} of $s$ on the episodes that are produced or that survive the filter, which changes with the input law and with selection (Theorem~\ref{thm:geometry-selection}, Remark~\ref{rem:instrument-vs-audit}).
  \item \textbf{Linear networks: Euclidean orthogonality suffices for the score function.} If $\varphi$ is the identity, the anchoring terms are $\psi_i = \nu\, c(x)(h_i - h_i^{(0)})^2$ with $\nu$ and $c$ common to all neurons and no parameter penalty, training is plain (full-batch or mini-batch) gradient descent or momentum from a zero buffer, and $W w = 0$, $U^\top w = 0$, then $w^\top (M(t) - M(0)) = 0$ for all $t$; again the score function is exactly invariant. The restriction on the optimizer is needed: Adam, which rescales coordinates, breaks the invariance (in a separate test, $200$ Adam steps with learning rate $10^{-3}$ on the linear companion network with the orthogonal shared probe give a drift $w^\top(M - M_0)$ of order $3 \times 10^{-6}$, against $2 \times 10^{-17}$ for gradient descent).
  \item \textbf{Nonlinear networks: orthogonality does not suffice.} Let $I := S \cup \{i : W_{:,i} \neq 0\}$, suppose $U_{i,:} = 0$ for $i \in I$, $\E_{\mathcal{D}}\|x\|^2 < \infty$ and $\varphi'$ bounded. For one gradient step of size $\gamma$ from $M$, the first-order change of the probe score at input $x'$ due to the filter path is
  \begin{equation}\label{eq:ntk-drift}
    \delta s(x') = -\gamma\, \E_{x \sim \mathcal{D}}\Big[\sum_{i \in I} w_i\, W_{:,i}^\top\, \nabla_z \ell_x \; \varphi'(m_i^\top x)\, \varphi'(m_i^\top x')\, x^\top x'\Big] + o(\gamma).
  \end{equation}
  If the kernels $k_i(x, x') := \varphi'(m_i^\top x)\varphi'(m_i^\top x')\, x^\top x'$, $i \in I$, are linearly independent in $L^2(\mathcal{D} \otimes \mathcal{D})$ (this fails, for instance, if a ReLU neuron in $I$ is dead, $k_i \equiv 0$), then $\delta s \equiv 0$ for every filter loss if and only if $w_i W_{:,i} = 0$ for all $i \in I$ (neuron-level disjointness). In the linear case all $k_i$ coincide and the condition collapses to $\sum_i w_i W_{:,i} = W w = 0$. This is the first-layer neural tangent kernel~\cite{jacot2018} restricted to the filter path.
  \item \textbf{A trainable chain-of-thought head recruits the probe's neurons.} If $W$ is trained jointly by gradient descent on~\eqref{eq:train-objective}, then along the gradient flow $\tfrac{d}{dt} W_{:,S} = -\E_x\big[\nabla_z \ell_x\; h_S(x)^\top\big]$, the non-centred cross moment between the filter-loss gradient and the activations of the probe's neurons. Neuron-disjointness is therefore not preserved: starting from $W_{:,S} = 0$, the head acquires weight on the probe's neurons, after which Part~1 no longer applies. The quantity to monitor is $\|W_{:,S}\|$ (or $\sum_{i \in S} |w_i|\, \|W_{:,i}\|$), not the coupling $\|Ww\|$: in the companion network, for a probe chosen after the fact orthogonal to the rows of $W(T)_{:,S}$ (legitimate since neither $W$ nor $M$ depends on $w$), $\|W(T) w\| = 5 \times 10^{-16}$ while $\|W(T)_{:,S}\| = 3.06$ and the probe score drifts by $3.4 \times 10^{-3}$. The coupling $\|Ww\|$ is the relevant quantity only for linear networks (Part~2).
\end{enumerate}
\end{theorem}

Part~1 is the elementary fact that parameters with identically zero gradient do not move; it is the mechanism used to mask gradients in gradient routing~\cite{cloud2024}. Part~4 is a direct computation.

\begin{proof}
\emph{Part 1.} $\partial \mathcal{L}/\partial m_i = \E_x\big[\big((W^\top \nabla_z \ell_x)_i + (U \nabla_y \ell_x)_i + \partial_1\psi_i(h_i(x), x)\big)\varphi'(m_i^\top x)\, x\big] + \nabla r_i(m_i)$, where $\nabla_z, \nabla_y$ are the gradients of $\ell_x$ in its two arguments. For $i \in S$ the first two terms vanish because $(W^\top v)_i = W_{:,i}^\top v = 0$ and $(U v)_i = U_{i,:} v = 0$. If $m_i(t') = m_i(0)$ for all $t' \le t$, then $h_i(x) = h_i^{(0)}(x)$ and the anchoring terms vanish, so every gradient with respect to $m_i$ up to time $t$ is zero; from a zero optimizer state the update of $m_i$ is then zero; induction on $t$ (or, for the gradient flow, uniqueness of the solution of the ODE restricted to $\{m_i = m_i(0),\ i \in S\}$, which is invariant) concludes. The cancellation is pointwise in $x$, so mini-batches and policy-dependent inputs are covered.

\emph{Part 2.} With $\varphi = \mathrm{id}$, $\nabla_M \mathcal{L} = \E_x\big[(W^\top \nabla_z\ell_x + U \nabla_y \ell_x + 2\nu\, c(x) (M - M_0) x)\, x^\top\big]$. Multiplying on the left by $w^\top$: the first term gives $(Ww)^\top \E[\nabla_z \ell_x x^\top] = 0$, the second $(U^\top w)^\top \E[\nabla_y\ell_x x^\top] = 0$, the third $2\nu\, w^\top (M - M_0)\E[c(x) x x^\top]$, which is zero if $w^\top(M - M_0) = 0$. The update is a linear function of past gradients for gradient descent and momentum, so induction over steps (or invariance of the linear subspace $\{w^\top (M - M_0) = 0\}$ under the flow) gives the claim; $s(x) = w^\top M x$ is then constant in $t$.

\emph{Part 3.} The filter-path gradient with respect to $m_i$ is $\E_x[W_{:,i}^\top \nabla_z \ell_x\, \varphi'(m_i^\top x)\, x]$, so one step changes $h_i(x') = \varphi(m_i^\top x')$ by $-\gamma\, \varphi'(m_i^\top x')\, x'^\top \E_x[W_{:,i}^\top\nabla_z\ell_x\,\varphi'(m_i^\top x)\, x] + o(\gamma)$ (for the ReLU, at inputs $x'$ off the hyperplane $m_i^\top x' = 0$, a $\mathcal{D}$-null set when $\mathcal{D}$ has a density); summing $w_i \times$ this over $i \in I$ gives~\eqref{eq:ntk-drift}. If $w_i W_{:,i} = 0$ for all $i$, the drift vanishes. Conversely, the filter loss can be chosen so that $x \mapsto \nabla_z \ell_x$ is any bounded measurable $a : \mathbb{R}^p \to \mathbb{R}^m$ (e.g.\ $\ell_x(z, y) = a(x)^\top z$); the integrability assumptions make each $k_i(\cdot, x')$ square-integrable. Vanishing of $\delta s(x')$ for all such $a$ and $\mathcal{D}$-a.e.\ $x'$ means, for each output coordinate $j$, $\sum_{i \in I} w_i W_{ji}\, k_i(\cdot, x') = 0$ in $L^2(\mathcal{D})$ for a.e.\ $x'$, hence $\sum_{i \in I} w_i W_{ji} k_i = 0$ in $L^2(\mathcal{D} \otimes \mathcal{D})$ by Fubini, and linear independence forces $w_i W_{ji} = 0$. When $\varphi = \mathrm{id}$, $k_i = x^\top x'$ for all $i$ and the condition is $\sum_i w_i W_{:,i} = Ww = 0$.

\emph{Part 4.} $\partial \mathcal{L}/\partial W = \E_x[\nabla_z \ell_x\, h(x)^\top]$, so under the flow $\dot W = -\E_x[\nabla_z\ell_x h(x)^\top]$; restricting to the columns in $S$ gives the stated equation. A probe $w \neq 0$ supported on $S$ and orthogonal, at time $t$, to the rows of $W(t)_{:,S}$ (such a $w$ exists whenever $|S| > \operatorname{rank} W(t)_{:,S}$) has $W(t) w = 0$ although $W(t)_{:,S} \neq 0$, which is why $\|Ww\|$ cannot detect recruitment.
\end{proof}

\begin{remark}[An unchanged instrument is not an unchanged audit]\label{rem:instrument-vs-audit}
Under Operator~A, the quantity of the static regime (R2) is audit survival among the cheating episodes that survive the filter, a property of the joint law of filter survival and probe score. Theorem~\ref{thm:geometry-training}(1) fixes the probe's score function, but training changes the filter's survival function $\beta_{\mathrm{loop}}(x) = 1 - \sigma(a^\top W h(x) - b)$, hence the selection weight of Operator~A and the law of $s$ among survivors; Theorem~\ref{thm:geometry-selection} must be re-applied at every step with the current $\beta_{\mathrm{loop}}$. A constructive example (two cheating subtypes, one seen mostly by the filter and one mostly by the probe; ReLU network; probe on neurons disjoint from the head; audit threshold at the 95th percentile of honest scores before training; script \texttt{v4\_contre\_exemple\_signe.py}) keeps the probe's rows identical to the bit and its area under the ROC curve at $0.691$, yet the audit's catch rate among filter survivors falls from $0.519$ to $0.356$ and joint evasion (passing both filter and audit) rises from $0.185$ to $0.492$: training makes more of the filter-visible subtype survive, which the probe does not see. On the companion network, by contrast, the same quantity moves favourably (from $0.359$ to $0.415$ for the disjoint ReLU probe), and in $48$ further runs ($12$ seeds, the companion task and a two-direction variant, linear and ReLU, $600$ steps) it never decreased by more than $10^{-3}$; the companion experiment therefore does not stress this risk.
\end{remark}

\begin{remark}[What the geometry says about practice]
These statements apply known tools to the filter-versus-audit setting; their content concerns dependence created by shared computation rather than the placement of the audit. First, the invariance of a probe's score function under training against the filter is a property of \emph{which neurons} the probe and the text head read, not of the angle between their readout vectors: in a ReLU layer, a probe orthogonal to the chain-of-thought readout but sharing its neurons is moved by the filter gradient for generic filter losses, whereas a probe on disjoint neurons keeps its score function for the whole run. Second, keeping the score function is not keeping the audit: the audit among filter survivors can still degrade by selection (Remark~\ref{rem:instrument-vs-audit}), so an auditor should measure audit survival weighted by filter survival, not only the probe's overall area under the curve. Third, the guarantee is fragile to training the text head: the head can recruit the probe's neurons, so an auditor relying on Part~1 must either freeze the head's access to those neurons or monitor $\|W_{:,S}\|$, which by Part~4 is observable during training; monitoring $\|Ww\|$ is not enough. These statements concern one hidden layer; in deeper networks, earlier shared layers couple all neurons and only the first-order formula~\eqref{eq:ntk-drift} (with the corresponding tangent kernel) remains. Empirical work on the evasion of latent-space monitors under optimization~\cite{bailey2024,gupta2025} studies related phenomena in real models; the present section gives exact statements only in a toy architecture.
\end{remark}

\paragraph{Numerical check.} The companion script \texttt{experience\_jouet\_decouplage\_sonde\_v1.py} verifies~\eqref{eq:stein-shift} by Monte Carlo ($4 \times 10^6$ draws; deviations from the prediction of $0.5$, $1.2$ and $1.1$ standard errors for aligned, $\Sigma$-orthogonal and random probes, with a crude standard error that is not centred and ignores that the prediction is estimated on the same sample; by one-dimensional quadrature the identity holds to $10^{-10}$), the counterexample of Remark~\ref{rem:uncorrelated}, and Theorem~\ref{thm:geometry-training} on a network with $d = 32$ hidden units trained for $1500$ full-batch gradient steps against a logistic filter on a two-dimensional text head. The episodes are fixed inputs: only the representation is trained, not the behaviour. Results (frozen-probe area under the ROC curve, cheating versus honest, before $\to$ after; filter catch rate on cheating episodes before $\to$ after):
\begin{itemize}
  \item neuron-disjoint probe, linear and ReLU: probe scores unchanged exactly (the change is $0$ to the bit), while the filter catch rate falls ($0.507 \to 0.316$ for ReLU). This is a check of the code, not a test that could have failed: the gradient rows of the probe's neurons are zero by construction;
  \item orthogonal probe sharing the text head's neurons: unchanged in the linear network (drift $1.8 \times 10^{-15}$, as Part~2 guarantees), moved but not evaded in the ReLU network (maximal score change $0.78$, shift of the mean cheater score $-0.058$, area under the curve $0.749 \to 0.750$);
  \item probe aligned with the text readout: evaded together with the filter (area under the curve $0.708 \to 0.657$ linear, $0.565 \to 0.435$ ReLU);
  \item a re-fitted probe keeps an area under the curve of about $0.83$--$0.86$ in all cases. This too follows from the construction: in the linear network $0.856$ is the Bayes area under the curve on the inputs, $\Phi(1.5/\sqrt{2}) = 0.8556$, since cheating is an input label and $h = Mx$ generically stays injective ($M$ of rank $p$); in the ReLU network the re-fitted area stays at its initial value ($0.832 \to 0.833$--$0.835$), below the Bayes value, because the map $h = \mathrm{ReLU}(Mx)$ with $d = 2p$ loses little information during training. It does not show that separability survives when episodes depend on the policy;
  \item with a trainable text head, the filter is evaded almost entirely through the head (catch rate $0.507 \to 0.0003$) and the head recruits the probe's neurons: $\|W_{:,S}\|$ grows from $0$ to $3.06$ while the coupling $\|Ww\|$ grows only from $0$ to $0.62$, with a small probe-score drift ($4.5 \times 10^{-3}$) over this horizon. (The companion figure still plots $\|Ww\|$; its next version should plot $\|W_{:,S}\|$.)
\end{itemize}
